\section{From the block tree to the map}\label{metric:section}

The mean block-distance comparison follows the general strategy for
subcritical graph classes~\cite[Section~5]{PanagiotouStuflerWeller2016}
and tree-like maps~\cite[preprint Section~6.8]{Stufler2020}.
Here Theorem~\ref{tree:paths} supplies the required uniformity under
only finite offspring variance. We also retain the exact corner-mass
calculation, as in~\cite[Section~9]{Stufler2026}, because the conclusion
specifies the degree probability measure.

Let $\mathsf T_N$ be the conditioned marked tree of
Proposition~\ref{block:finite-law}, where $N=2n+1$. Its map has the
required marginal law of $M_n$. Write $\pi(w)$ for the pole of a tree
node $w$. For a strict ancestor with $k$ children, block mark $B$, and
selected child $i$, put
\[
 g(k,i,B)=d_B\big(\text{pole of }B,\text{vertex of corner }i\big),
 \qquad \beta=\E_{\rm sp}g.
\]
The expectation is under the spine law \eqref{tree:spine-law}, and
$d_B$ uses all block edges. Since a block has $k/2$ edges,
$0\le g\le k$. Therefore
\begin{equation}\label{metric:beta}
 0<\beta<\infty.
\end{equation}
Finiteness follows from $\E_{\rm sp}k=\E\xi^2<\infty$.
For positivity, a one-link block and its corner at the endpoint
opposite the pole have positive spine probability, and their distance
is one.

Let $G(w)$ be the sum of these $g$ labels over the strict ancestors of
$w$. By Theorem~\ref{tree:paths} and the coding estimates in
Section~\ref{tree:section}, with $\nu=\sigma^2/2$,
\begin{equation}\label{metric:uniform-poles}
 \max_{w\in\mathsf T_N}
 \left|G(w)-\beta H(w)\right|=\op(\sqrt N),\qquad
 \max_{w\in\mathsf T_N}
 \left|G(w)-\frac{\beta}{\nu}W(w)\right|=\op(\sqrt N).
\end{equation}

\subsection{Uniform comparison of distances}

\begin{lemma}[Distances between poles]\label{metric:poles}
Let $w_*$ be the last common ancestor of two tree nodes $w,w'$.
If $\Delta$ is the maximum offspring number, then
\[
 \left|d_{M_n}\big(\pi(w),\pi(w')\big)
       -\big(G(w)+G(w')-2G(w_*)\big)\right|\le2\Delta.
\]
The map $\pi$ from tree nodes to map vertices is onto.
\end{lemma}

\begin{proof}
Every inserted map meets the preceding part at its pole only.
An excursion from a block into an attached submap must return to the
same vertex, so attachments cannot shorten distances within a block.
Moving from an ancestor pole to a descendant pole therefore adds
the distances $g$ along their chain of blocks.

If one of $w,w'$ is their last common ancestor, this proves the
displayed comparison with zero error. Otherwise let $x,x'$ be the two
selected corner vertices in the block of $w_*$. The route through
the pole of that block has length
$d_B(x,\pi(w_*))+d_B(\pi(w_*),x')$ within the block, whereas the
shortest route uses $d_B(x,x')$. Their difference lies between zero
and twice the diameter of $B$, and that diameter is at most its
number of edges, hence at most $\Delta$. Outside the common block
the distances add exactly as above. The same argument allows
coincident poles and loop blocks.

Every vertex of the nonempty map is incident to at least one edge,
and hence to a corner in one of its blocks. That corner is a child
slot, whose child's pole is the given vertex. Thus every map vertex
is the pole of a nonroot node.
\end{proof}

Index the nodes by $j=0,\ldots,N-1$ in preorder. Combining
Lemma~\ref{metric:poles}, \eqref{metric:uniform-poles}, and
Lemma~\ref{tree:lca}, and using $\Delta=\op(\sqrt N)$, gives
\begin{equation}\label{metric:walk-comparison}
 \max_{0\le j\le j'<N}
 \left|d_{M_n}\big(\pi(j),\pi(j')\big)
  -\frac{\beta}{\nu}
       \left(W_j+W_{j'}-2\min_{j\le r\le j'}W_r\right)\right|
   =\op(\sqrt N).
\end{equation}
In particular, the comparison is uniform over every map vertex,
including vertices in the largest blocks.

\subsection{The degree measure}

\begin{lemma}[Corner mass]\label{metric:measure}
The uniform probability measure on the $N-1=2n$ nonroot tree nodes,
pushed forward by $\pi$, is exactly $\mu_n$.
\end{lemma}

\begin{proof}
Each nonroot node occupies one corner of its parent block, and every
block corner is occupied once, including the corners whose inserts
are trivial leaves. At a map vertex $v$, summing these corner counts
over the blocks gives $\deg_{M_n}(v)$: each incident dart belongs to
one block, and each block has one corner per incident dart. A loop
contributes two darts and two corners at its vertex.
Division by the total $2n$ proves the assertion.
\end{proof}

If $\bar\mu_n$ is instead the pushforward of uniform probability on
all $N$ nodes, then
\begin{equation}\label{metric:mass-error}
 \bar\mu_n=\frac{N-1}{N}\mu_n+\frac1N\delta_{\pi(0)},
 \qquad
 \|\bar\mu_n-\mu_n\|_{\rm TV}\le\frac1N.
\end{equation}
The measure identity is thus independent of the sizes and shapes of
the individual blocks.

\subsection{A correspondence estimate}

For completeness we state the elementary estimate used to pass to
metric probability spaces. A \emph{correspondence} between $X$ and
$Y$ is a relation $\mathcal R\subseteq X\times Y$ projecting onto both
spaces. Its distortion is
\[
 \operatorname{dis}(\mathcal R)
 =\sup_{\substack{(x,y),(x',y')\in\mathcal R}}
       |d_X(x,x')-d_Y(y,y')|.
\]
The Gromov--Hausdorff--Prokhorov distance is the infimum, over common
isometric embeddings, of the maximum of the Hausdorff distance
between the embedded spaces and the Prokhorov distance between their
probability measures.

\begin{lemma}[Metric and mass comparison]\label{metric:ghp}
Suppose a correspondence between compact metric probability spaces
has distortion at most $\varepsilon>0$, and a coupling of their
probability measures is supported on the correspondence. Their
Gromov--Hausdorff--Prokhorov distance is at most $\varepsilon$.
If one marginal of that coupling differs from the desired measure
by total variation at most $\eta$, the bound is
$\varepsilon+\eta$.
\end{lemma}

\begin{proof}
Join corresponding points by links of length $\varepsilon$, and
use the induced shortest-chain distance on the disjoint union.
This preserves the distances within each original space.
Indeed, a segment leaving one space and returning to it can be
replaced by a segment inside that space: the two linking lengths
are $2\varepsilon$, whereas the difference of the distances between
the corresponding endpoints is at most $\varepsilon$.
Every point is within $\varepsilon$ of the other space, giving the
Hausdorff bound. The coupling places every pair at distance at most
$\varepsilon$, which gives the Prokhorov bound directly from its
defining neighborhood inequalities. Changing one marginal by
total variation $\eta$ increases these inequalities, and hence the
Prokhorov bound, by at most $\eta$.
\end{proof}

For a nonnegative continuous function $f$ on $[0,1]$ with
$f(0)=f(1)=0$, define $d_f$ by the same formula as $d_e$, and denote
its quotient metric probability space by $\mathcal T_f$.
The interval-minimum formula gives the triangle inequality for $d_f$.
The quotient projection is continuous, since
$d_f(s,t)$ is at most twice the oscillation of $f$ on the interval
between $s$ and $t$. Thus $\mathcal T_f$ is compact.
Relating equal-time representatives in $\mathcal T_f$ and
$\mathcal T_{\tilde f}$ gives a correspondence of distortion at most
$4\|f-\tilde f\|_\infty$ and a coupling from one uniform time.
Lemma~\ref{metric:ghp} therefore proves continuity of
\[
 f\longmapsto \mathcal T_f
\]
from the uniform norm to the Gromov--Hausdorff--Prokhorov topology.

\subsection{Proof of the main theorem}

\begin{proof}[Proof of Theorem~\ref{thm:main}]
Define
\begin{equation}\label{metric:constant}
 c(q)=\frac{\nu}{\sqrt2\,\beta\sigma}
     =\frac{\sigma}{2\sqrt2\,\beta}.
\end{equation}
Equations~\eqref{block:moments} and \eqref{metric:beta} show that this
constant is deterministic, positive, and finite for every fixed $q>4$.

Let $f_N$ interpolate $W_j/(\sigma\sqrt N)$ at $j/N$ for
$0\le j<N$, and set $f_N(1)=0$. It is nonnegative and has zero
endpoints. Replacing the exploration walk's last value
$W_N=-1$ by zero changes its normalized interpolation by at most
$1/(\sigma\sqrt N)$. Proposition~\ref{tree:excursion} consequently gives
\[
 f_N\Longrightarrow e
 \qquad\text{in }C([0,1]).
\]
For grid points $j/N,j'/N$, the comparison
\eqref{metric:walk-comparison} and the identity $N=2n+1$ imply
\[
 \max_{j,j'<N}
 \left|
  \frac{c(q)}{\sqrt n}\,
     d_{M_n}\big(\pi(j),\pi(j')\big)
  -\sqrt{\frac{N}{2n}}\,
     d_{f_N}(j/N,j'/N)
 \right|=\op(1).
\]
Since $\sqrt{N/(2n)}\to1$ and $\|f_N\|_\infty$ is tight, that
deterministic factor may be replaced by one.

For $s\in[0,1]$, put $j_N(s)=\min\{\lfloor Ns\rfloor,N-1\}$ and
write $[s]$ for its class in $\mathcal T_{f_N}$. The relation
\[
 \mathcal R_N=\{(\pi(j_N(s)),[s]):s\in[0,1]\}
\]
is a correspondence by Lemma~\ref{metric:poles}.
Rounding times changes a $d_{f_N}$ distance by at most
$4\omega_{f_N}(1/N)$, where $\omega$ is the uniform modulus of
continuity. Continuous-path tightness makes this error tend to zero
in probability. The correspondence therefore has distortion
$\op(1)$ for the rescaled map metric.

A Lebesgue-uniform time couples the probability measure of
$\mathcal T_{f_N}$ to $\bar\mu_n$ exactly, since each of the $N$
preorder cells has length $1/N$. By
Lemma~\ref{metric:ghp} and \eqref{metric:mass-error}, the
Gromov--Hausdorff--Prokhorov distance from the rescaled map with
measure $\mu_n$ to $\mathcal T_{f_N}$ tends to zero in probability.
The continuity established above and $f_N\Longrightarrow e$ prove
the claimed convergence to $(\mathcal T_e,d_e,\mu_e)$.

Proposition~\ref{block:finite-law} applies to every $n\ge1$, and
the conditioned-tree results hold through all the corresponding
odd sizes. The proof uses only the map marginal and all-edge
block distances, so its conclusion is precisely the unrooted,
undecorated metric probability limit stated in the theorem.
\end{proof}
